\section*{Bloque A. Fundamentos teóricos y conceptuales}

\noindent\textit{Preguntas 1 a 30. Evalúan la comprensión analítica del valor del dinero en el tiempo, regímenes de interés, TEA y ecuación de Fisher, anualidades y perpetuidades, estados financieros, ratios de liquidez, solvencia y rentabilidad (DuPont), flujos de caja libres y criterios de evaluación de proyectos (VAN, TIR, TIRM, Payback, CAE y WACC).}

\q{El principio del 'Valor del Dinero en el Tiempo' establece fundamentalmente que una unidad monetaria hoy vale más que una unidad monetaria en el futuro debido a:}
\begin{opts}
    \item existencia exclusiva de intermediarios bancarios con encaje legal positivo.
    \item mandato normativo discrecional fijado anualmente por el Banco Central.
    \item preferencia por consumo presente, inflación esperada y costo de oportunidad del capital.
    \item volatilidad cambiaria esperada respecto a monedas extranjeras líquidas.
\end{opts}

\q{En el régimen de interés compuesto frente al de interés simple:}
\begin{opts}
    \item los intereses devengados se capitalizan y generan nuevos intereses en periodos sucesivos.
    \item la tasa de crecimiento del capital acumulado sigue una trayectoria estrictamente lineal.
    \item en interés compuesto no se cumple la propiedad de aditividad temporal del flujo de caja.
    \item los intereses generados en cada periodo se retiran sin integrarse al capital inicial.
\end{opts}

\q{La Tasa Efectiva Anual (TEA) se distingue de la Tasa Nominal Anual (TNA) en que la TEA:}
\begin{opts}
    \item es fijada por mandato administrativo anual del Ministerio de Economía y Finanzas.
    \item ignora la frecuencia de capitalización subperiódica de intereses dentro del año fiscal.
    \item mide el costo efectivo anual incorporando capitalización ($TEA = (1 + TNA/m)^m - 1$).
    \item es siempre estrictamente menor que la TNA cuando la capitalización es mensual o diaria.
\end{opts}

\q{La Ecuación de Fisher exacta que relaciona la tasa de interés nominal ($i$), la tasa de interés real ($r$) y la tasa de inflación esperada ($\pi$) se expresa formalmente como:}
\begin{opts}
    \item $1 + r = \frac{1 + i}{\pi}$
    \item $1 + i = (1 + r)(1 + \pi)$
    \item $1 + i = \frac{1 + r}{1 + \pi}$
    \item $r = (1 + i)(1 + \pi) - 1$
\end{opts}

\q{El Valor Presente ($VP$) de una corriente perpetua de flujos de caja constantes e iguales a $C$, que se cobran a partir del final del primer periodo con una tasa de descuento constante $r$, es exactamente:}
\begin{opts}
    \item $VP = \frac{C}{(1 + r)^n}$
    \item $VP = \frac{C}{r - g}$
    \item $VP = C \times r$
    \item $VP = \frac{C}{r}$
\end{opts}

\q{El modelo de Gordon y Shapiro para el Valor Presente de una perpetuidad creciente a una tasa constante $g$ anual a perpetuidad ($C_1$ a cobrar en $t=1$, con $r > g$) establece que:}
\begin{opts}
    \item $VP = C_1 \times (r - g)$
    \item $VP = \frac{C_1}{g}$
    \item $VP = \frac{C_1}{r - g}$
    \item $VP = \frac{C_1}{r + g}$
\end{opts}

\q{La relación matemática entre el Valor Presente de una anualidad anticipada ($VP_{ant}$) y el Valor Presente de una anualidad vencida u ordinaria ($VP_{venc}$) con idéntico número de cuotas y tasa de interés $r$ es:}
\begin{opts}
    \item $VP_{ant} = VP_{venc} \times (1 - r)$
    \item $VP_{ant} = VP_{venc} \times (1 + r)$
    \item $VP_{ant} = VP_{venc} \times r$
    \item $VP_{ant} = VP_{venc} / (1 + r)$
\end{opts}

\q{La Curva de Rendimiento (Yield Curve) que representa la estructura temporal de las tasas de interés suele tener pendiente positiva en condiciones macroeconómicas normales porque:}
\begin{opts}
    \item los inversionistas institucionales prefieren inmovilizar capital a largo plazo por liquidez.
    \item el Banco Central impone por norma una estructura temporal de tasas estrictamente creciente.
    \item los plazos largos conllevan mayor riesgo de tasa e inflación, exigiendo prima de liquidez.
    \item los emisores soberanos de largo plazo carecen de pasivos contingentes en su balance general.
\end{opts}

\q{En el Balance General de una empresa, la ecuación contable básica que expresa la igualdad fundamental de fuentes y usos de recursos es:}
\begin{opts}
    \item Activo = Pasivo + Patrimonio Neto
    \item Pasivo = Activo + Patrimonio Neto
    \item Patrimonio = Activo + Pasivo Neto
    \item Activo = Pasivo - Patrimonio Neto
\end{opts}

\q{La 'Prueba Ácida' (o ratio de liquidez severa) se diferencia de la 'Razón Corriente' en que la prueba ácida:}
\begin{opts}
    \item excluye los inventarios del activo corriente por ser el rubro realizable menos líquido.
    \item deduce los préstamos bancarios de largo plazo del pasivo corriente.
    \item multiplica las cuentas por cobrar comerciales por el margen bruto anual.
    \item suma los activos fijos depreciables al disponible neto de caja y bancos.
\end{opts}

\q{El 'Capital de Trabajo Neto' (Net Working Capital) de una empresa prestadora de servicios se define operativamente como:}
\begin{opts}
    \item el flujo de caja neto operativo generado durante el ejercicio contable.
    \item la diferencia entre Activo Corriente y Pasivo Corriente ($AC - PC$).
    \item la suma de las utilidades retenidas más las reservas legales del capital.
    \item el total de activos fijos tangibles netos de depreciación acumulada.
\end{opts}

\q{El ratio de Apalancamiento Financiero Deuda/Patrimonio ($D/E$) mide:}
\begin{opts}
    \item el financiamiento ajeno frente a fondos propios, reflejando el riesgo financiero.
    \item el número de veces que el capital de trabajo neto rota respecto a las ventas anuales totales.
    \item la rentabilidad operativa neta obtenida por cada sol invertido en activos intangibles.
    \item la proporción de activos corrientes que han sido liquidados en efectivo al cierre fiscal.
\end{opts}

\q{El ratio de Cobertura de Intereses (Times Interest Earned) se calcula como:}
\begin{opts}
    \item $\frac{EBITDA}{Activo\ Corriente}$
    \item $\frac{EBIT}{Gastos\ Financieros}$
    \item $\frac{Utilidad\ Neta}{Pasivo\ Total}$
    \item $\frac{Ventas}{Gastos\ Financieros}$
\end{opts}

\q{La diferencia conceptual entre la Rentabilidad sobre los Activos ($ROA$) y la Rentabilidad sobre el Patrimonio ($ROE$) radica en que:}
\begin{opts}
    \item el ROE mide exclusivamente el impacto de la inflación sectorial sobre el capital de trabajo de la empresa.
    \item el ROA se aplica únicamente a bancos comerciales mientras el ROE mide liquidez en empresas reales.
    \item el ROA supera siempre al ROE cuando la firma contrata deuda bancaria con escudo fiscal positivo.
    \item el ROA evalúa la eficiencia de los activos totales y el ROE el retorno sobre el patrimonio propio.
\end{opts}

\q{La descomposición DuPont del $ROE$ desagrega la rentabilidad sobre el patrimonio en el producto de tres motores financieros clave:}
\begin{opts}
    \item Margen bruto operativo $\times$ Razón corriente $\times$ Inflación esperada anual.
    \item Margen Neto $\times$ Rotación de Activos $\times$ Multiplicador de Capital ($A/P$).
    \item Costo de ventas $\times$ Tasa impositiva efectiva $\times$ Cobertura de deuda ($EBIT$).
    \item Flujo de caja libre $\times$ Tasa de interés nominal $\times$ Costo promedio de capital.
\end{opts}

\q{El 'Ciclo de Conversión del Efectivo' (Cash Conversion Cycle) mide el tiempo transcurrido desde que la empresa invierte efectivo en insumos hasta que recupera el efectivo por cobranza, calculándose como:}
\begin{opts}
    \item Días de inventario + Días de cobro - Días de pago.
    \item Rotación de inventarios + Días de pago / 365.
    \item Días de cobro - Días de inventario + Días de pago.
    \item Días de cobro + Periodo promedio de pago a proveedores.
\end{opts}

\q{Al proyectar el Flujo de Caja Libre del Proyecto (FCLP o Free Cash Flow to Firm - FCFF):}
\begin{opts}
    \item suma la utilidad neta contable después de intereses más el total de pasivos no corrientes.
    \item descuenta íntegramente las cuotas de amortización de préstamos y los gastos financieros.
    \item excluye el servicio de la deuda (intereses y amortización), midiendo el valor operativo puro.
    \item incluye como egresos relevantes los costos hundidos devengados en ejercicios anteriores.
\end{opts}

\q{En la preparación del flujo de caja de un proyecto de inversión, un 'Costo Hundido' (Sunk Cost), como el dinero ya gastado en un estudio geológico previo no reembolsable:}
\begin{opts}
    \item debe excluirse del flujo incremental por ser un gasto histórico irrecuperable.
    \item debe registrarse como inversión inicial negativa en el periodo $t=0$.
    \item debe actualizarse al costo de oportunidad e incluirse como flujo operativo.
    \item debe depreciarse tributariamente de forma lineal durante diez años.
\end{opts}

\q{El criterio del Valor Actual Neto (VAN / NPV) para la toma de decisiones de inversión establece que un proyecto financieramente independiente debe aceptarse si:}
\begin{opts}
    \item la relación Beneficio/Costo resulta estrictamente menor que la unidad ($B/C < 1$).
    \item $VAN < 0$, indicando una menor salida de capital respecto al costo ponderado.
    \item el valor actual de los flujos iguala al costo inicial ($VAN = 0$) en sector privado.
    \item $VAN > 0$, pues los flujos descontados superan la inversión y crean riqueza.
\end{opts}

\q{Una propiedad matemática y financiera fundamental que hace que el criterio del VAN sea superior a otros indicadores es la 'Aditividad de Valor', la cual significa que:}
\begin{opts}
    \item el VAN de todo proyecto aumenta monótonamente al subir el costo de capital.
    \item para proyectos independientes, se cumple $VAN(A+B) = VAN(A) + VAN(B)$.
    \item la TIR conjunta de un portafolio equivale a la suma de las TIR individuales.
    \item el VAN resulta insensible a los flujos de caja de horizontes lejanos.
\end{opts}

\q{La Tasa Interna de Retorno (TIR / IRR) se define rigurosamente como:}
\begin{opts}
    \item la tasa de descuento que anula el Valor Actual Neto ($VAN(TIR) = 0$).
    \item la tasa de interés pasiva fijada por el banco comercial más grande del mercado.
    \item el cociente contable promedio entre la utilidad bruta y los activos totales.
    \item la tasa que maximiza la suma aritmética simple de flujos de caja nominales.
\end{opts}

\q{Una limitación técnica crítica de la Tasa Interna de Retorno frente al VAN radica en el supuesto implícito de reinversión de los flujos intermedios, según el cual:}
\begin{opts}
    \item la TIR asume reinversión a tasa cero mientras el VAN supone ausencia total de reinversión.
    \item el VAN asume reinversión al costo de capital ($WACC$), mientras la TIR asume reinversión a la propia TIR.
    \item ambos métodos suponen reinversión obligatoria a la tasa libre de riesgo de los bonos soberanos.
    \item el VAN asume reinversión a la tasa de inflación; la TIR asume reinversión al costo de oportunidad del capital.
\end{opts}

\q{En proyectos con flujos de caja no convencionales (donde ocurren múltiples cambios de signo en los flujos a lo largo del tiempo, por ejemplo por costos de cierre de una mina o planta de tratamiento), la Regla de los Signos de Descartes implica que:}
\begin{opts}
    \item el proyecto carece de Valor Actual Neto definido para tasas reales positivas.
    \item el proyecto debe rechazarse automáticamente sin calcular su Valor Actual Neto.
    \item el VAN es estrictamente cóncavo y creciente respecto a la tasa de descuento relevante.
    \item pueden existir múltiples TIR (tantas como cambios de signo), invalidando ese criterio.
\end{opts}

\q{Al comparar dos proyectos de inversión mutuamente excluyentes con diferente escala o diferente perfil temporal de flujos (donde las curvas de VAN se intersectan en la llamada Tasa de Fisher), el criterio que debe prevalecer siempre es:}
\begin{opts}
    \item el proyecto que requiera la mayor inversión inicial en infraestructura física.
    \item el Periodo de Recuperación simple más corto para minimizar el riesgo operativo.
    \item la Tasa Interna de Retorno (TIR), eligiendo el proyecto con mayor porcentaje de rentabilidad.
    \item el Valor Actual Neto (VAN) al costo de capital, por maximizar el valor absoluto.
\end{opts}

\q{La Tasa Interna de Retorno Modificada (TIRM / MIRR) supera las limitaciones de la TIR tradicional porque:}
\begin{opts}
    \item prescinde de conocer la vida útil del proyecto y la estructura de financiamiento de la firma.
    \item actualiza egresos al financiamiento y capitaliza ingresos al costo de capital, dando tasa única.
    \item elimina del flujo los desembolsos negativos iniciales para asegurar convergencia numérica.
    \item transforma automáticamente cualquier secuencia de flujos para que arroje un VAN positivo.
\end{opts}

\q{El método del Periodo de Recuperación del Capital (Payback simple) presenta dos deficiencias teóricas fundamentales para la toma de decisiones óptimas:}
\begin{opts}
    \item sobrestima sistemáticamente el capital de trabajo y la depreciación.
    \item ignora el valor del dinero en el tiempo y los flujos posteriores al recupero.
    \item exige cálculo matricial iterativo y depende de la deuda contable.
    \item asume la liquidación forzosa de activos fijos al cumplir el plazo.
\end{opts}

\q{El Costo Anual Equivalente (CAE / EAC) es la herramienta financiera idónea para:}
\begin{opts}
    \item comparar proyectos mutuamente excluyentes con distintas vidas útiles y costos.
    \item medir el impacto de la inflación acumulada sobre los contratos de construcción civil.
    \item evaluar la solvencia de corto plazo de una firma ante choques de capital de trabajo.
    \item determinar el precio de emisión de acciones preferentes en el mercado secundario.
\end{opts}

\q{La fórmula clásica del Costo Promedio Ponderado de Capital (WACC / CPPC) calculada sobre la estructura de capital a valores de mercado es:}
\begin{opts}
    \item $WACC = K_e (1 - T) \left(\frac{E}{V}\right) + K_d \left(\frac{D}{V}\right)$
    \item $WACC = K_e \left(\frac{E}{V}\right) + K_d (1 - T) \left(\frac{D}{V}\right)$
    \item $WACC = K_e \left(\frac{E}{V}\right) + \frac{K_d}{(1 - T)} \left(\frac{D}{V}\right)$
    \item $WACC = K_e \left(\frac{E}{V}\right) + K_d \left(\frac{D}{V}\right)$
\end{opts}

\q{En el modelo Capital Asset Pricing Model (CAPM) para el cálculo del costo del capital propio ($K_e$), la ecuación representativa es $K_e = R_f + \beta (R_m - R_f)$. El parámetro Beta ($\beta$) mide:}
\begin{opts}
    \item el riesgo diversificable o específico propio de la gestión interna de la firma.
    \item la volatilidad idiosincrásica del tipo de cambio frente al oro monetario.
    \item el riesgo sistemático no diversificable respecto al portafolio de mercado.
    \item la rentabilidad nominal libre de riesgo de los bonos soberanos del tesoro.
\end{opts}

\q{La Proposición I de Modigliani y Miller (1958) en un mundo ideal de mercados de capitales perfectos (sin impuestos corporativos, sin costos de quiebra y con información simétrica) establece que:}
\begin{opts}
    \item las empresas deben financiarse exclusivamente con utilidades retenidas.
    \item el valor de mercado crece linealmente al aumentar el ratio de deuda.
    \item el valor de mercado de la firma es totalmente independiente de su estructura de capital.
    \item el costo promedio ponderado de capital equivale a la tasa libre de riesgo.
\end{opts}

